\documentclass[a4paper,11pt]{ctexart}

% 引入必要的宏包
\usepackage[utf8]{inputenc}
\usepackage[T1]{fontenc}
\usepackage{geometry}
\usepackage{amsmath, amssymb, amsfonts} % 数学公式支持
\usepackage{bm} % 粗体数学符号
\usepackage{hyperref} % 超链接支持
\usepackage{url}

% additional packages
\usepackage{tabularx}
\usepackage{booktabs}
\usepackage{multirow}
\usepackage{graphicx}

% 页面边距设置
\geometry{left=2.5cm, right=2.5cm, top=2.5cm, bottom=2.5cm}

% 标题设置
\title{\textbf{Mergenet 写作提纲}}
\author{}
\date{}

\begin{document}

\maketitle

\begin{abstract}
transformers在建模语言方面表现卓越，但在内存和时间复杂度方面效率低下。一种可能的解决方案是缩短序列长度。（Dynamic Token Pooling，Hourglass Transformer）分词就是最有效的缩短序列长度，让模型精准理解文本内容的手段，但是当前的分词器面临着一些局限性。现有的 LLM 严重依赖静态分词器（BPE，UniGram，SentencePiece），导致了跨领域适应性差、无法处理 Out-of-Distribution (OOD) 文本，且在处理信息密度不均文本时计算效率低下（由于这些分词器基于频率而非语义，导致模型在处理“高频无意义词”和“低频高义词”时分配了同样的算力，同时 Rho-1 已经向我们证明“并非所有 Token 都重要”）。

面对这样的问题，端到端学习是一种出路。端到端学习模型（MegaByte，SpaceByte，MambaByte，BLT，HNet）尝试移除分词器，但面临“硬切分”导致的语义断裂、梯度无法回传优化分词策略，以及推理时无法动态修正前期错误的局限。

\textbf{Mergenet 核心创新：}提出一种全局可微的层级架构。
\begin{itemize}
    \item 引入 DTEM (Decoupled Token Embedding Merging) 机制，在 LoE 阶段通过可学习的度量动态聚类 Byte 序列。
    \item 利用 Perceiver-style Cross-Attention （Perceiver）将聚合后的语义“概念”压缩为固定长度的 Latent Vectors，实现高压缩率。
    \item LoD连续滑动解码结合重建损失，确保生成过程的语义对齐。
\end{itemize}

\textbf{结果预告：}在保持 byte-level 灵活性的同时，实现了与 word-level 模型相当的推理速度，并在代码/多语言等 OOV 密集场景下显著优于 BPE 模型，比别的端到端模型如BLT更优。
\end{abstract}

\section{Introduction}

\subsection{The Tokenization Bottleneck}
阐述 BPE/SentencePiece 的“静态假设”在面对新颖词汇、代码、DNA 序列时的失效。
指出 Tokenization 本质上是对算力的预分配，静态分配导致了算力与信息密度的不匹配（高频无义词占用过多算力）。

\subsection{From Hard Boundaries to Soft Concepts}
批评现有 Byte-level 方法（MegaByte, BLT）虽然去除了词表，但依赖启发式或预测性的“硬边界”，破坏了语义的连续性，且难以端到端优化切分逻辑。分词不应是预处理，而应是模型“理解”过程的一部分（Concept Formation）。

\subsection{Mergenet Proposal}
概述 Mergenet 架构：Local Encoder (Byte $\to$ "Concept, not only sub-word") $\to$ Latent Encoder (Concept Processing, like formal GPT) $\to$ Decoder (Concept $\to$ Byte)。
全链路可微，基于内容的动态压缩，LaE 层的 KV Cache 复用同普通GPT模型。

\section{Related Work}

\textbf{Subword-based LLMs:} 简述 BPE 的局限性。

\textbf{Byte-level Modeling:} MegaByte，SpaceByte，MambaByte，BLT，HNet。

\textbf{Dynamic Tokenization \& Boundary Learning:} 对比 BLT 和 H-Net。

\textbf{Critical Differentiation:} 指出它们是从左到右预测边界（不可逆，无梯度回传），而 Mergenet 是基于全局/局部注意力矩阵的聚类（可微，软对齐）。

\textbf{Vision Transformers \& Token Merging:} 提及 TokenLearner, Perceiver, 和 ViT 中的 Token Merging (ToMe)，阐述 Mergenet 如何将视觉领域的“空间聚合”迁移至 NLP 的“时序语义聚合”。

\subsection{与长上下文语言模型（LCLM）综述脉络的对应}
近期综述 Liu et al.\ （arXiv:2503.17407）将长上下文研究按数据、架构、工作流、基础设施、评测与行为分析等维度组织（RQ1--RQ3）。
\textbf{MergeNet 的主定位在 RQ1 的架构侧：}在\textbf{表示层}通过可微合并（DTEM）与 Perceiver bottleneck 缩短有效序列长度、降低 LaE 注意力成本，属于``序列抽象 / 高效注意力''分支，而非推理期 KV 压缩（\S5）、亦非外挂式 Prompt 压缩或 RAG 流水线（工作流 \S4）---写作时需在 Related Work 中显式划界。
LaE 可选用线性注意力、窗口注意力或混合结构，与综述 \S3 中次线性/混合架构对应；连续质心 $c_i$ 与弹性网格偏置则与位置编码、外推及\textit{有效上下文长度}（如 RULER、lost-in-the-middle）的讨论相关，实验上应联合报告压缩率 $\lambda$、推理窗 $W_{infer}$ 与下游长任务，避免仅宣称原始 byte 窗长。
\textbf{数据与训练（综述 \S2）}可作为后续迭代：长文过滤、长度配比与 GrowLength 式课程是否适配 byte 噪声与高压缩率，需在扩展实验中单独验证（详见 \texttt{MergeNet/LCLM综述与MergeNet对齐.md} 与 \texttt{optimization.tex} 补充清单）。

\section{Methods}

\subsection{IO-Aware Rank-1 Bias Injection via Dimension Augmentation}
传统的相对位置编码或窗口偏置通常需要显式实例化一个偏置矩阵 $\mathbf{B} \in \mathbb{R}^{N \times N}$，这会导致 $O(N^2)$ 的显存开销，并破坏 FlashAttention 等 IO 感知算子的融合机制。针对本架构中广泛存在的“秩-1 (Rank-1)”偏置结构（如基于对数窗口的 $\beta_j = \log S_j$），我们提出了一种维度增广策略，将加法偏置转化为高维空间的矩阵乘法。

我们构建增广的 Query 和 Key 矩阵 $\hat{\mathbf{Q}}, \hat{\mathbf{K}} \in \mathbb{R}^{N \times (d+1)}$：
\begin{equation}
    \hat{\mathbf{Q}} = [\mathbf{Q} \mid \sqrt{d} \cdot \mathbf{1}_N], \quad \hat{\mathbf{K}} = [\mathbf{K} \mid \boldsymbol{\beta}]
\end{equation}
其中 $\mathbf{1}_N$ 为全 1 列向量，$\boldsymbol{\beta}$ 为偏置向量。对增广矩阵执行标准的 Scaled Dot-Product Attention，其结果在数学上严格等价于注入偏置：
\begin{equation}
    \frac{\hat{\mathbf{Q}}\hat{\mathbf{K}}^\top}{\sqrt{d}} = \frac{\mathbf{Q}\mathbf{K}^\top + (\sqrt{d}\cdot\mathbf{1}_N)\boldsymbol{\beta}^\top}{\sqrt{d}} = \text{Score}(\mathbf{Q}, \mathbf{K}) + \mathbf{1}_N \boldsymbol{\beta}^\top
\end{equation}
该方法将内存复杂度严格控制在 $O(N)$，且无需修改底层的 CUDA Attention Kernel，完美兼容 FlashAttention 的分块计算逻辑。

\subsection{Continuous Geometric Alignment \& Elastic Grid}
由于 MergeNet 摒弃了离散的 Token 索引，我们需要在连续空间中重建位置感知能力。

\paragraph{Physical Centroid Calculation (LoE):}
我们将 Token 的“位置”重定义为其所聚合的 Raw Bytes 的物理质心。位置索引 $c_i$ 不再是整数，而是由源矩阵 $\mathbf{S}$ 和质量向量 $\mathbf{v}$ 决定的动态函数：
\begin{equation}
  c_i = \frac{\sum_k S_{i,k} \cdot k}{\mathbf{v}_i + \epsilon}
\end{equation}
这意味着位置信息本身是可微的。模型可以通过调整 $\mathbf{S}$（即改变聚类范围）来微调 Concept 在语义空间中的逻辑位置，实现“位置”的端到端学习。

\paragraph{Elastic Grid with Soft Bias (LoD):}
在解码阶段，为了解决 Concept 与 Byte 的非刚性对应问题，我们摒弃了硬性的边界映射，转而采用基于距离的软网格偏置（Soft Grid Bias）。解码器维护一个基于压缩率 $\lambda$ 的弹性网格，偏置定义为：
\begin{equation}
  \text{Bias}_{grid}(t, j) = - \gamma \cdot \left| \frac{t}{\lambda} - j \right|
\end{equation}
配合滑动窗口 $\text{SlidingWin}_{grid}$，模型能够关注到当前 Byte 附近的 Latent Concepts。这不仅实现了推理时的平滑对齐，还允许通过梯度下降自动学习最佳的关注跨度（Attention Span）。

\subsection{Parallelized Dynamic Merging \& Storage Optimization}
\paragraph{Parallel Top-K Merging:}
原本 DTEM 模拟将模型中 top-k 对最相关的token取出来并进行可微合并需要循环 k 次，高复杂度低并行度，效率低。针对 DTEM 中动态聚类涉及的高昂循环 Top-K 操作，我们引入了基于阈值的并行化处理（Threshold-based Parallel Top-K, inspired by Su et al.）。通过将串行的循环选择转化为可并行计算的阈值掩码操作，解决了动态分词中的计算瓶颈，大幅提升了 LoE 阶段的吞吐量。（\url{https://spaces.ac.cn/archives/10373}）

\paragraph{Quasi-Linear Complexity with Unfold Storage:}
采用 SlidingWinAttn，同时创新性地使用unfold将带状 size 矩阵和 tokenmap 存储到狭长矩阵中，保证了计算和存储高效，保证了内存访问的连续性。

\subsection{System-Level Hardware Adaptation}
\textbf{Hardware Alignment Padding:} 为适配 GPU Tensor Cores 的计算特性，我们将增广维度 $d+1$ 自动填充至最近的 8 的倍数 ($d_{\text{phys}}$)。填充位全零初始化，确保在利用硬件加速特性的同时不引入计算噪声。

\textbf{Zero-Copy Inference Mechanism:} 在推理阶段，我们引入持久化 KV Cache 队列 $\mathcal{Q} \in \mathbb{R}^{B \times (2W_{infer}+1) \times d}$。通过预分配显存并采用原地修改策略，直接将计算出的偏置项写入预留通道，实现了前向传播过程中的零内存分配和零数据拷贝。

\section{Experiments}
评测与报告协议对齐 LCLM 综述 \S6--\S7 思路：除 LongBench 与 common sense 外，应披露 $\lambda$、$W_{infer}$ 及 LaE 算子配置，并区分合成任务与真实长任务（详见 \texttt{MergeNet/LCLM综述与MergeNet对齐.md}）。
\begin{table*}[t]
\caption{Performance comparison on language modeling and zero-shot common-sense reasoning. Avg. denotes the average of accuracy without perplexity.}
\centering
\setlength{\tabcolsep}{1.0mm}
\resizebox{1.\linewidth}{!}{
    \begin{tabular}{l|c|cc|cccCCccccccc}
    \toprule
    \multirow{2}*{\textbf{Model}} & Params. 
    & Wiki. & LMB. & LMB. & PIQA & Hella. & Wino. & ARC-e & ARC-c & SIQA & BoolQ & COPA & OBQA & Avg. \\
    & (M) & ppl $\downarrow$ & ppl $\downarrow$ & acc $\uparrow$ & acc $\uparrow$ & acc\_n $\uparrow$ & acc $\uparrow$ & acc $\uparrow$ & acc\_n $\uparrow$ & acc $\uparrow$ & acc $\uparrow$ & acc $\uparrow$ & acc\_n $\uparrow$ &  \\
    \midrule
    \multicolumn{15}{c}{BPE Tokenizer} \\
    \midrule
    Transformer++ & 370M & 89.13 & 26.27 & 36.89 & 66.27 & 41.89 & 52.88 & 59.51 & 28.41 & 85.90 & 60.31 & 74.00 & 33.80 & 53.98  \\
    GSA & 380M & 288.11 & 69395.11 & 31.00 & 55.33 & 25.08 & 48.62 & 35.23 & 22.44 & 37.10 & 43.18 & 59.00 & 24.40 & 38.14 \\
    GLA & 370M & 31.55 & 36.34 & 31.73 & 64.80 & 38.14 & 49.57 & 56.19 & 25.85 & 81.60 & 58.84 & 68.00 & 31.80 & 50.65 \\
    DeltaNet & 370M & 29.99 & 36.89 & 31.83 & 66.32 & 38.23 & 55.41 & 56.31 & 27.47 & 81.20 & 58.62 & 71.00 & 34.00 & 52.04 \\
    GDN & 512M & 25.67 & 27.67 & 34.50 & 67.25 & 42.22 & 49.17 & 59.89 & 28.24 & 85.40 & 61.35 & 70.00 & 35.20 & 53.32 \\
    GDN-H & 460M  & 25.62 & 26.07 & 36.29 & 67.25 & 41.83 & 53.43 & 60.77 & 29.95 & 85.70 & 57.28 & 70.00 & 33.60 & 53.61 \\
    \midrule
    \multicolumn{15}{c}{Byte-level Tokenizer} \\
    \midrule
    Transformer++ & 370M & 875.93 & 34.55 & 31.40 & 62.95 & 36.29  & 51.70 & 51.98 & 26.45 & 83.90 & 41.62 & 62.00 & 29.40 & 47.76 \\
    GSA & 380M & 71.44 & 107.94 & 16.20 & 59.19 & 30.88 & 50.43 & 43.90 & 22.27 & 77.80 & 44.95 & 56.00 & 28.80 & 43.04\\
    GLA & 370M & 48.21 & 47.04 & 26.26 & 62.57 & 34.03 & 50.43 & 49.20 & 25.34 & 78.30 & 40.67 & 60.00 & 30.40 & 45.72 \\
    DeltaNet & 370M & 44.06 & 40.68 & 27.05 & 61.97 & 34.50 & 50.43  & 51.98 & 24.49 & 80.90 & 41.31 & 62.00 & 32.20 & 46.68 \\
    GDN & 512M & 36.96 & 29.22 & 30.04 & 64.85 & 37.76 & 51.22 & 54.88 & 26.54 & 83.50 & 39.72 & 64.00 & 33.20 & 48.57 \\
    GDN-H & 460M  & 37.03 & 29.61 & 31.71 & 63.98 & 37.56 & 51.14 & 54.29 & 23.34 & 84.30 & 46.03 & 64.00 & 32.40 & 48.87 \\
    BLT & 380M & 44.30 & 42.35 & 29.19 & 61.48 & 34.21 & 50.20 & 46.76 & 24.06 & 82.90 & 38.78 & 58.00 & 31.00 & 45.66 \\
    \bottomrule
    \end{tabular}
}
\label{tab:350m_results}
\end{table*}
\begin{table*}[t]
\caption{Performance comparison on language modeling and LongBench.}
\centering
\setlength{\tabcolsep}{1.0mm}
\resizebox{\linewidth}{!}{
    \begin{tabular}{l|c|ccc|ccc|ccc|ccc|cc|c}
    \toprule
    \multirow{2}*{\textbf{Model}} & Params. & \multicolumn{3}{c|}{Single-Doc QA} & \multicolumn{3}{c|}{Multi-Doc QA} & \multicolumn{3}{c|}{Summarization} & \multicolumn{3}{c|}{Few-shot} & \multicolumn{2}{c|}{Code} & Avg \\
    & (M) & NQA & QQA & MFQ & HQA & 2WM & Mus & GvR & QMS & MNs & TRC & TQA & SSM & LCC & RBP & \\
    \midrule
    \multicolumn{17}{c}{BPE Tokenizer} \\
    \midrule
    Transformer++ & 370M & 0.68 & 1.89 & 4.47 & 0.95 & 1.48 & 0.30 & 5.39 & 2.58 & 8.70 & 2.50 & 3.43 & 2.30 & 12.79 & 5.05 & 3.75 \\
    GSA & 380M & 0.56 & 1.05 & 2.76 & 0.84 & 0.90 & 0.71 & 6.17 & 10.94 & 8.44 & 0.00 & 4.02 & 3.09 & 11.98 & 9.63 & 4.36 \\
    GLA & 370M & 1.71 & 3.42 & 9.84 & 3.77 & 6.97 & 2.51 & 6.99 & 14.19 & 9.32 & 10.50 & 14.65 & 6.42 & 17.42 & 15.49 & 8.80 \\
    DeltaNet & 370M & 1.79 & 3.57 & 11.34 & 4.32 & 6.41 & 2.60 & 5.51 & 15.18 & 10.06 & 13.00 & 13.32 & 7.46 & 16.87 & 20.53 & 9.43 \\
    GDN & 512M & 2.17 & 3.74 & 13.16 & 4.17 & 7.81 & 3.10 & 6.81 & 17.48 & 10.50 & 23.00 & 17.96 & 19.01 & 22.39 & 23.40 & 12.48 \\
    GDN-H & 460M & 1.81 & 4.29 & 13.69 & 4.59 & 8.89 & 2.89 & 6.83 & 15.21 & 11.14 & 27.00 & 19.79 & 16.57 & 19.93 & 19.42 & 12.29 \\
    \bottomrule
    \end{tabular}
}
\label{tab:350m_results_longbench}
\end{table*}

\subsection{Datasets}

\subsection{Setups}

\subsection{Implement Deatils}

\end{document}
